\section{Introduction}
\label{sec:introduction}

The electroweak sector contains a dimensionless mass ratio,
\begin{equation}
  \sPole = 1-\frac{M_{W,\rm pole}^2}{M_{Z,\rm pole}^2},
  \label{eq:pole-angle-intro}
\end{equation}
which is fixed experimentally by the pole positions of the charged and neutral weak vector bosons.  In the Standard Model at tree level the same ratio is written as an angular datum in coupling space, while the Higgs vacuum expectation value fixes the common radial scale of the vector masses.  This separation between projective and radial data gives the working language of the present paper.

The Rivero--de Vries observation starts from a mass operator built from Poincare Casimir data and produces a dimensionless two-branch equation.  In the normalized notation used below, the equation is
\begin{equation}
  x^2+Jx-J=0,
  \qquad J=s(s+1).
  \label{eq:intro-secular}
\end{equation}
The positive branch gives
\begin{equation}
  \sdV = 1-\frac{x_+(3/4)}{x_+(2)}=0.22310132\ldots .
  \label{eq:intro-devries}
\end{equation}
The original note emphasized the proximity of this number to the mass-shell weak angle and motivated the operator through Casimir invariants and high-spin behavior \cite{Rivero2006}.  The present manuscript reorganizes the point as a pole-spectrum statement and asks what physics can turn Eq.~\eqref{eq:intro-secular} into an electroweak calculation.

The construction has three layers.  The first layer is algebraic: a two-channel secular equation has roots whose ratio is fixed by the Casimir variable.  The second layer is electroweak: the W/Z mass ratio is a projective datum along the broken-to-unbroken ray of the gauge-Higgs system.  The third layer is dynamical: strings, branes, Kaluza--Klein boundary conditions, or singular \(G_2\) geometries may generate the secular equation or its branch structure.

The status of the argument is deliberately stratified.  The algebraic calculation is explicit.  The numerical comparison is reproducible.  The pole-mass scheme must be source-audited.  The assignment of \(J=3/4\) to the charged W comparison and \(J=2\) to the neutral Z comparison is the central open analytical calculation unless it is derived in Sec.~\ref{sec:ew-embedding}.  The negative branch is retained because it may encode the scalar/order-parameter side of electroweak symmetry breaking.

The paper is organized as follows.  Section~\ref{sec:casimir} gives the secular operator and its roots.  Section~\ref{sec:pole} defines the pole observable.  Section~\ref{sec:ew-embedding} formulates the projective electroweak interpretation.  Section~\ref{sec:negative} records the negative branch and its possible Higgs-scale role.  Section~\ref{sec:string-kk-g2} states the string/Kaluza--Klein/\(G_2\) completion problem.  Sections~\ref{sec:flavour} and \ref{sec:global-form} delimit the flavor and global-gauge-group issues.  Section~\ref{sec:status-tests} gives the status of the analytical program and falsifiable tests.

\paragraph*{Section status.}  This introduction states the program and labels the central unproved step.
